\documentclass[10pt,a4paper]{amsart}
\usepackage{amsmath,amssymb,amsthm,graphicx}
\usepackage{verbatim}
\newenvironment{reference}{\comment}{\endcomment}
%\newenvironment{reference}{}{}
\newenvironment{slogan}{\comment}{\endcomment}
\newenvironment{history}{\comment}{\endcomment}

% For commutative diagrams we use Xy-pic
\usepackage[all]{xy}

% We use 2cell for 2-commutative diagrams.
\xyoption{2cell}
% Theorem environments.
%
\theoremstyle{plain}
\newtheorem{theorem}[subsection]{Theorem}
\newtheorem{proposition}[subsection]{Proposition}
\newtheorem{lemma}[subsection]{Lemma}

\theoremstyle{definition}
\newtheorem{definition}[subsection]{Definition}
\newtheorem{example}[subsection]{Example}
\newtheorem{exercise}[subsection]{Exercise}
\newtheorem{situation}[subsection]{Situation}

\theoremstyle{remark}
\newtheorem{remark}[subsection]{Remark}
\newtheorem{remarks}[subsection]{Remarks}

\numberwithin{equation}{subsection}

% Macros
%
\def\lim{\mathop{\mathrm{lim}}\nolimits}
\def\colim{\mathop{\mathrm{colim}}\nolimits}
\def\Spec{\mathop{\mathrm{Spec}}}
\def\Hom{\mathop{\mathrm{Hom}}\nolimits}
\def\Ext{\mathop{\mathrm{Ext}}\nolimits}
\def\SheafHom{\mathop{\mathcal{H}\!\mathit{om}}\nolimits}
\def\SheafExt{\mathop{\mathcal{E}\!\mathit{xt}}\nolimits}
\def\Sch{\mathit{Sch}}
\def\Mor{\mathop{\mathrm{Mor}}\nolimits}
\def\Ob{\mathop{\mathrm{Ob}}\nolimits}
\def\Sh{\mathop{\mathit{Sh}}\nolimits}
\def\NL{\mathop{N\!L}\nolimits}
\def\CH{\mathop{\mathrm{CH}}\nolimits}
\def\proetale{{pro\text{-}\acute{e}tale}}
\def\etale{{\acute{e}tale}}
\def\QCoh{\mathit{QCoh}}
\def\Ker{\mathop{\mathrm{Ker}}}
\def\Im{\mathop{\mathrm{Im}}}
\def\Coker{\mathop{\mathrm{Coker}}}
\def\Coim{\mathop{\mathrm{Coim}}}

% Boxtimes
%
\DeclareMathSymbol{\boxtimes}{\mathbin}{AMSa}{"02}

%
% Macros for moduli stacks/spaces
%
\def\QCohstack{\mathcal{QC}\!\mathit{oh}}
\def\Cohstack{\mathcal{C}\!\mathit{oh}}
\def\Spacesstack{\mathcal{S}\!\mathit{paces}}
\def\Quotfunctor{\mathrm{Quot}}
\def\Hilbfunctor{\mathrm{Hilb}}
\def\Curvesstack{\mathcal{C}\!\mathit{urves}}
\def\Polarizedstack{\mathcal{P}\!\mathit{olarized}}
\def\Complexesstack{\mathcal{C}\!\mathit{omplexes}}
% \Pic is the operator that assigns to X its picard group, usage \Pic(X)
% \Picardstack_{X/B} denotes the Picard stack of X over B
% \Picardfunctor_{X/B} denotes the Picard functor of X over B
\def\Pic{\mathop{\mathrm{Pic}}\nolimits}
\def\Picardstack{\mathcal{P}\!\mathit{ic}}
\def\Picardfunctor{\mathrm{Pic}}
\def\Deformationcategory{\mathcal{D}\!\mathit{ef}}
\usepackage{cleveref}
\crefname{originalreference}{original reference}{original references}
\Crefname{originalreference}{Original reference}{Original references}
\makeatletter
\newlabel{lemma-F-zero-module-map}{{083L}{source}{}{sourcepointer-1}{}}
\newlabel{lemma-F-zero-module-map@cref}{{[originalreference][0][]083L}{[1][0][]source}{}{}{}}
\makeatother
\crefname{originalsection}{original section}{original sections}
\Crefname{originalsection}{Original section}{Original sections}
\makeatletter
\newlabel{more-morphisms-section-finite-type-finite-presentation}{{05GX}{source}{}{sourcepointer-2}{}}
\newlabel{more-morphisms-section-finite-type-finite-presentation@cref}{{[originalsection][0][]05GX}{[1][0][]source}{}{}{}}
\makeatother
\crefname{originalreference}{original reference}{original references}
\Crefname{originalreference}{Original reference}{Original references}
\makeatletter
\newlabel{proposition-finite-presentation-flat-at-point}{{05ID}{source}{}{sourcepointer-3}{}}
\newlabel{proposition-finite-presentation-flat-at-point@cref}{{[originalreference][0][]05ID}{[1][0][]source}{}{}{}}
\makeatother
\crefname{originalreference}{original reference}{original references}
\Crefname{originalreference}{Original reference}{Original references}
\makeatletter
\newlabel{situation-somewhat-closed}{{0CWA}{source}{}{sourcepointer-4}{}}
\newlabel{situation-somewhat-closed@cref}{{[originalreference][0][]0CWA}{[1][0][]source}{}{}{}}
\makeatother
\title{A map that vanishes on a set of points vanishes at every point of its closure}
\author{The Stacks Project Authors}
\date{stacks-project:spaces-flat.tex:local-snapshot-20261012; GNU FDL 1.2 or later}
\begin{document}
\maketitle
\noindent\textit{Editorial note.} All delivered mathematical text is the author's own, copied from the pinned Stacks Project chapter file spaces-flat.tex; no statement or proof step was rewritten, reordered or removed.\par\medskip
\noindent\textit{Editorial note.} The counted proof is the author's complete argument: an etale local reduction to a projective chart with the closure condition preserved along the open morphism, and an algebra conclusion through the ideal of the map\par\medskip
\noindent\textit{Editorial note.} The label lemma-F-zero-somewhat-closed-points names the counted statement, and the proof environment immediately following it in the source is delivered verbatim as its proof.\par\medskip
\noindent\textit{Editorial note.} The counted statement is carried by the author's own lemma environment in the chapter source.\par\medskip
\noindent\textit{Allowed standard prerequisites (editorial).}\begin{itemize}
\item The situation of a map of quasi-coherent modules over the base with flat target, and the pullback of the map to the fibres over points of the base
\item Etale locality of the vanishing of the pullback of a module map on a fibre, and the existence of an etale neighbourhood on which the source becomes projective over the base with a prescribed point of the fibre covered
\item The ideal of a map between a module and a projective module, and the description of its vanishing locus as the set of points where the pulled back map is zero
\end{itemize}\medskip
\ifdefined\hypertarget\hypertarget{sourcepointer-1}{}\fi
\noindent\textit{Source reference.} Original reference 083L: sources/stacks-spaces-flat/pinned-original-source.tex, source lines 1627--1627. Prerequisite pointer invoked at proof line(s) 1685; it names a standard fact of the same chapter that the counted argument uses, and every step of the counted proof built on it is retained verbatim in the delivered body.\par\medskip
\ifdefined\hypertarget\hypertarget{sourcepointer-2}{}\fi
\noindent\textit{Source reference.} Original section 05GX: sources/stacks-chow/tags.txt, source lines 7076--7076. Prerequisite pointer invoked at proof line(s) 1660; it names a standard fact of a sibling chapter that the counted argument uses, and every step of the counted proof built on it is retained verbatim in the delivered body.\par\medskip
\ifdefined\hypertarget\hypertarget{sourcepointer-3}{}\fi
\noindent\textit{Source reference.} Original reference 05ID: sources/stacks-spaces-flat/pinned-original-source.tex, source lines 924--924. Prerequisite pointer invoked at proof line(s) 1674; it names a standard fact of the same chapter that the counted argument uses, and every step of the counted proof built on it is retained verbatim in the delivered body.\par\medskip
\ifdefined\hypertarget\hypertarget{sourcepointer-4}{}\fi
\noindent\textit{Source reference.} Original reference 0CWA: sources/stacks-spaces-flat/pinned-original-source.tex, source lines 1560--1560. Prerequisite pointer invoked at proof line(s) 1649; it names a standard fact of the same chapter that the counted argument uses, and every step of the counted proof built on it is retained verbatim in the delivered body.\par\medskip
\begin{lemma}
\label{lemma-F-zero-somewhat-closed-points}
In Situation \ref{situation-somewhat-closed}.
Let $T \subset S$ be a subset. Let $s \in S$ be in the closure of $T$.
For $t \in T$, let $u_t$ be the pullback of $u$ to $X_t$
and let $u_s$ be the pullback of $u$ to $X_s$.
If $X$ is locally of finite presentation over $S$,
$\mathcal{G}$ is of finite presentation\footnote{It would
suffice if $X$ is locally of finite type over $S$
and $\mathcal{G}$ is finitely presented relative to $S$,
but this notion hasn't yet been defined in the setting
of algebraic spaces. The definition for schemes is
given in More on Morphisms, Section
\ref{more-morphisms-section-finite-type-finite-presentation}.}, and
$u_t = 0$ for all $t \in T$, then $u_s = 0$.
\end{lemma}

\begin{proof}
To check whether $u_s$ is zero, is \'etale local on the fibre $X_s$.
Hence we may pick a point $x \in |X_s| \subset |X|$ and check
in an \'etale neighbourhood. Choose
$$
\xymatrix{
(X, x) \ar[d] & (X', x') \ar[l]^g \ar[d] \\
(S, s) & (S', s') \ar[l]
}
$$
as in Proposition \ref{proposition-finite-presentation-flat-at-point}.
Let $T' \subset S'$ be the inverse image of $T$. Observe that
$s'$ is in the closure of $T'$ because $S' \to S$ is open.
Hence we reduce to the algebra problem described in the
next paragraph.

\medskip\noindent
We have an $R$-module map $u : M \to N$ such that $N$ is projective
as an $R$-module and such that
$u_t : M \otimes_R \kappa(t) \to N \otimes_R \kappa(t)$
is zero for each $t \in T$. Problem: show that $u_s = 0$.
Let $I \subset R$ be the ideal defined in Lemma \ref{lemma-F-zero-module-map}.
Then $I$ maps to zero in $\kappa(t)$ for all $t \in T$.
Hence $T \subset V(I)$. Since $s$ is in the closure of $T$
we see that $s \in V(I)$. Hence $u_s = 0$.
\end{proof}

\end{document}
